%% ═══════════════════════════════════════════════════════════════════════════
\section{A Worked Example, End to End}
\label{app:worked}
%% ═══════════════════════════════════════════════════════════════════════════

Section~\ref{sec:math_framework} gives the model as equations and Section~\ref{sec:Results}
gives it as scores on a campaign. Neither shows a single number travelling through the whole
chain, which is what a reader needs to check that the equations mean what the prose
says they mean. This appendix carries one node through every stage: the six blocks into a
local urgency, that urgency along a three-node chain, the chain into a criticality ranking,
and the result against a declaration.

Every figure below is produced by the production code rather than computed by hand, using
the same functions the campaign ran, so the example cannot drift away from the system it
illustrates. The scenario is illustrative and is not drawn from either campaign.

\subsection{Stage 1: six blocks into a local urgency}

The node is a deep-tier supplier whose next milestone falls in nine days against a lead time
distributed $\mathcal{N}(11, 3)$, whose buffers are a third consumed, whose OEE stands at
$0.80$, and for which no node-level emissions data exists. Table~\ref{tab:worked_blocks}
gives the resulting block values and the survival factor each contributes to
Equation~\ref{eq:ur_local}.

\begin{table}[htbp]
\centering
\caption[Worked example: the six blocks]{The six blocks and their survival factors, all
weights at $\omega_m = 1$}
\label{tab:worked_blocks}
\small
\begin{tabularx}{\linewidth}{@{}l c c X@{}}
\toprule
\textbf{Block} & \textbf{$u_m$} & \textbf{$(1-u_m)^{\omega_m}$} & \textbf{Reading} \\
\midrule
$u_{\text{time}}$ & $0.62$ & $0.3800$ & Probability of missing the nine-day deadline \\
$u_{\text{cap}}$  & $0.35$ & $0.6500$ & Buffers a third consumed \\
$u_{\text{perf}}$ & $0.20$ & $0.8000$ & $1 - \text{OEE}$, with OEE at $0.80$ \\
$u_{\text{risk}}$ & $0.15$ & $0.8500$ & Expected unavailability over the horizon \\
$u_{\text{cost}}$ & $0.28$ & $0.7200$ & Operating cost drifting above nominal \\
$u_{\text{co2}}$  & absent & $1.0000$ & $\omega_m = 0$: the factor is neutral, not healthy \\
\midrule
\multicolumn{2}{@{}l}{Product of survivals} & $0.120931$ & \\
\multicolumn{2}{@{}l}{$U_{r,\text{local}} = 1 - \text{product}$} & $\mathbf{0.8791}$ & \\
\bottomrule
\end{tabularx}
\end{table}

The comparison that matters is with the aggregator the model rejects. A weighted average of
the same five available blocks gives $0.3200$. The two answers differ by more than half the
scale on identical inputs, and they differ because they answer different questions: the
average asks how stressed this node is on average, the noisy-OR asks how likely it is that
at least one mechanism ruptures. Only the second is the quantity a coordinator needs.

The gap widens at the boundary. Let the deadline pass, so $u_{\text{time}} = 1$ and the node
cannot become any later. The noisy-OR returns exactly $1.0000$, because the factor
$(1-1) = 0$ collapses the product whatever the other blocks read. The weighted average
returns $0.3960$, and would rank this node below a node that is mildly stressed on all six.
That is the compensatory failure of Section~\ref{sec:model_ur} in a single pair of numbers.

\subsection{Stage 2: propagation along a three-node chain}

Place the node as the deepest tier of a chain $s_2 \to s_1 \to \text{customer}$, with
upstream amplification $\beta_{s_2 s_1} = 0.70$ and $\beta_{s_1 c} = 0.50$. The two
downstream nodes are locally calm, at $U_{r,\text{local}} = 0.30$ and $0.10$.
Equation~\ref{eq:ur_prop} applied in topological order gives
\begin{align*}
    U_{r,s_2} &= 0.8791, \\
    U_{r,s_1} &= 1 - (1 - 0.30)\bigl(1 - 0.70 \times 0.8791\bigr) = 0.7307, \\
    U_{r,c}   &= 1 - (1 - 0.10)\bigl(1 - 0.50 \times 0.7307\bigr) = 0.4288 .
\end{align*}
\begin{figure}[htbp]
\centering
\begin{tikzpicture}[x=1cm, y=1cm, font=\scriptsize,
    nd/.style = {bloc, text width=2.1cm, minimum height=1.05cm}]

\node[nd, fill=altenOchrePale!55!white, draw=altenOchreDark, line width=1pt] (s2) at (1.6, 0)
     {$s_2$\\[2pt]local $0.879$\\\textbf{prop. $0.879$}};
\node[nd] (s1) at (6.0, 0) {$s_1$\\[2pt]local $0.300$\\\textbf{prop. $0.731$}};
\node[nd, blocFort] (c) at (10.4, 0) {customer\\[2pt]local $0.100$\\\textbf{prop. $0.429$}};

\draw[fleche, line width=1.1pt] (s2) -- node[etiquette, above] {$\beta = 0.70$} (s1);
\draw[fleche, line width=1.1pt] (s1) -- node[etiquette, above] {$\beta = 0.50$} (c);

\node[etiquette, anchor=north, text width=3.4cm] at (1.6, -0.85)
     {six blocks, one absent\\(Table~\ref{tab:worked_blocks})};
\node[etiquette, anchor=north, text width=3.4cm] at (6.0, -0.85)
     {locally calm, inherits\\$0.43$ of its risk};
\node[etiquette, anchor=north, text width=3.6cm] at (10.4, -0.85)
     {its own indicators say\\nothing is wrong};

\end{tikzpicture}
\caption[The worked chain]{The three-node chain of Stage~2. Each box carries the node's own
local urgency above and its propagated value below. The customer's local reading of $0.10$
becomes $0.43$ once two tiers of upstream risk arrive, which is the whole purpose of the
propagation.}
\label{fig:worked_chain}
\end{figure}

A node that is locally calm at $0.10$ therefore carries a propagated urgency of $0.43$,
entirely inherited from two tiers upstream. This is the whole purpose of the propagation: the
customer's own indicators say nothing is wrong, and something is.

The log-survival twin of Proposition~\ref{prop:logsurvival} can be checked on the same
figures. Equation~\ref{eq:logsurvival} gives $\ell_c = 0.5601$ at the customer, and
$1 - e^{-0.5601} = 0.4288$, which is $U_{r,c}$ to every digit shown. The twin is the same
recurrence in another coordinate, exactly as the proposition states.

\subsection{Stage 3: which node is the most critical}

The criticality index of Section~\ref{sec:criticality} shocks each active node to
$U_{r,\text{local}} = 1$ in turn and measures the rise at the customer. On this chain:

\begin{table}[htbp]
\centering
\caption[Worked example: criticality]{Effect of the worst local shock, measured at the
customer}
\label{tab:worked_crit}
\small
\begin{tabularx}{\linewidth}{@{}l c c c X@{}}
\toprule
\textbf{Shocked node} & \textbf{$U_{r,\text{local}}$} & \textbf{$U_{r,c}$ after} &
\textbf{$\Delta U_{r,c}$} & \textbf{Rank} \\
\midrule
$s_1$ & $0.30 \to 1.00$ & $0.5500$ & $+0.1212$ & 1st \\
$s_2$ & $0.88 \to 1.00$ & $0.4555$ & $+0.0267$ & 2nd \\
\bottomrule
\end{tabularx}
\end{table}

The result is worth pausing on, because it is counter-intuitive and it is correct. The most
urgent node, $s_2$ at $0.88$, is the \emph{least} critical of the two. Two effects combine.
Its risk is already almost entirely realised in the baseline, so forcing it to $1$ adds
little that the customer was not already exposed to. And what it adds is damped twice, once
by $\beta_{s_2 s_1} = 0.70$ and again by $\beta_{s_1 c} = 0.50$, while $s_1$ sits behind a
single arc. Urgency and criticality are therefore different questions, and a coordinator who
sorts the network by urgency alone will spend attention on the node that has least left to
lose.

Proposition~\ref{prop:saturation} can be seen on the same chain. Set every local urgency to
$1.0$, the state the campaign reaches on nine turns of nineteen. The baseline customer
urgency is then $1$, and shocking either node gives $\Delta U_{r,c} = +0.0000$ exactly. Both
nodes tie at zero, the ranking carries no information, and the log-survival key of
Proposition~\ref{prop:order} is what separates them again.

\subsection{Stage 4: the declaration, and the adequacy it produces}

Finally, the deep supplier's own coordinator fills the weekly questionnaire.
Table~\ref{tab:worked_adequacy} scores two declarations against the computed
$U_{r} = 0.8791$, using the default $\lambda_{\text{under}} = 2.25$,
$\lambda_{\text{over}} = 1.0$ and $\alpha = 0.88$.

\begin{table}[htbp]
\centering
\caption[Worked example: adequacy]{Two declarations against the same computed urgency}
\label{tab:worked_adequacy}
\small
\begin{tabularx}{\linewidth}{@{}c c c c c X@{}}
\toprule
\textbf{$U_d$} & \textbf{$F$} & \textbf{$H$} & \textbf{penalty} & \textbf{$A$} &
\textbf{Reading} \\
\midrule
$0.83$ & $0.0000$ & $0.0491$ & $0.1585$ & $83.6$ & A lucid coordinator, slightly optimistic \\
$0.40$ & $0.0000$ & $0.4791$ & $1.1774$ & $22.7$ & Hidden risk of nearly half the scale \\
\bottomrule
\end{tabularx}
\end{table}

The second row is the situation the whole instrument exists to surface. The node is at
$0.88$ and is reported at $0.40$; nothing in the declaration is a lie, and the coordinator
may sincerely believe it, which is Assumption~\ref{as:sincerity}. The adequacy score falls to $22.7$ and the hidden risk
$H = 0.479$ is what a supervising view would rank on. Note also what the first row shows: at
$U_d = 0.83$ against $U_r = 0.88$ the score is $83.6$ and not $95$, because
Proposition~\ref{prop:adequacy}(iv) makes the curve steepest near zero. Small
mis-declarations are penalised hard on purpose, since they are the ones still cheap to
correct.
